\chapter{Visualization}
\label{chap:viz}

\section{The challenge}

The supraverse is a 4D manifold with non-trivial topology: a foam-like structure
of nested universe-bubbles connected by Cartasis hypersurfaces. This is hard to
visualize because 4D structures do not render to 2D screens directly; the foam
spans many orders of magnitude (subatomic to cosmic); universes sit inside black
holes inside universes inside black holes; time evolution within each universe
runs orthogonally to time evolution within others; and the supraverse as a whole
is static (a block universe) but locally evolves. Standard scientific
visualization techniques do not handle this well. New approaches are needed.

\section{The gravity-scaled conformal mapping approach}

\subsection{Core idea}

Apply a conformal transformation to the supraverse metric,
\[
  g_{\mu\nu}(x) \;\to\; \Omega^2(x)\, g_{\mu\nu}(x),
\]
where $\Omega(x)$ is a function of local gravitational field strength,
specifically
\[
  \Omega(x) \sim \bigl(G\,\rho(x)\bigr)^{-1/2}
\]
(the inverse square root of local mass--energy density, modulated by
gravitational coupling).

\subsection{What this does}

\begin{itemize}
  \item High-density regions (black-hole interiors, Cartasis membranes)
    \textbf{scale up} in the conformal map.
  \item Low-density regions (intergalactic voids, supraverse vacuum)
    \textbf{scale down}.
  \item The conformal transformation \textbf{preserves causal structure} (light
    cones map to light cones).
  \item Distances are distorted (this is a visualization tool, not a metric for
    physics).
\end{itemize}

\subsection{Visual result}

The supraverse, mapped this way, looks like bubbles of comparable visual size
regardless of physical scale; bounce membranes appear as prominent surfaces;
nested structure means zooming into any high-density region reveals more
universes; the appearance is self-similar and fractal-like; and chirality-coded
coloring distinguishes matter universes from antimatter universes.

\subsection{Relationship to Penrose diagrams}

Penrose diagrams are standard tools for visualizing general-relativistic
spacetimes. They use conformal transformations to compress infinite spacetime
onto finite diagrams while preserving causal structure. Penrose's Conformal
Cyclic Cosmology uses related techniques to glue cosmological cycles
together~\citep{penrose2010}. Our gravity-scaled mapping is a variant choice: use
\emph{gravity} as the conformal factor rather than choices made to fit spacetime
onto specific paper shapes. This emphasizes physically meaningful features
(high-density structures, bounces) rather than coordinate-system features
(asymptotic regions).

\section{Visualization targets}

\subsection{Single universe with nested black holes}

The simplest visualization. Take our universe and show internal structure with
galaxies and matter distributions. Within each galaxy, scale up to show black
holes; within each black hole, scale up to show the Cartasis membrane; inside the
membrane (geometrically inside but topologically connected to a separate region),
show the baby universe. This makes the Tardis effect immediately apparent:
outside scale small (Schwarzschild radius), inside scale large (baby cosmological
extent).

\subsection{Family tree of lineages}

Show our universe's ancestors going up the lineage (parent, grandparent,
great-grandparent) and our descendants going down (our internal black holes' baby
universes). Use alternating colors for matter/antimatter chirality across
generations. This makes the tree structure explicit and visualizes the
alternating-chirality prediction.

\subsection{Supraverse population}

Zoom out to show many independent lineages distributed across the supraverse---equal
numbers of matter-led and antimatter-led lineages, with vacuum between them. What that
view \emph{looks} like depends decisively on a single computed quantity, the bounce
instanton action $I$ (Chapter~\ref{chap:supraverse}, Section~\ref{sec:instanton}), and
this is worth stating carefully because the most intuitive picture is the wrong one.

\paragraph{The tempting picture (a packed polygonal foam) is not ours.} If OG
nucleation were efficient---$I<I_{\mathrm{crit}}\simeq1.4$, i.e.\ a near-Planckian
seed---the OGUs would nucleate faster than de Sitter expansion dilutes them, percolate,
and tile the void into a \emph{packed polygonal foam}: a soap-froth of curved-polygon
cells sharing membranes, coarsening over time. This is the visually rich, lava-lamp
image one naturally reaches for, and it would make a beautiful rendering. But the
computation puts $I\sim10^{84}\gg I_{\mathrm{crit}}$ (Section~\ref{sec:instanton}), so
this regime \emph{does not obtain}. We deliberately do \emph{not} show a packed-foam
figure here: it would depict a supraverse that the framework's own numbers exclude.

\paragraph{The actual picture is dilute.} For $I\gg I_{\mathrm{crit}}$ the void
inflates between rare nucleations, so OGUs are isolated, round, and astronomically far
apart---the dilute scatter of \Cref{fig:dilute} below, with the separation-to-diameter
ratio of $\sim10^{2.7\times10^{83}}$ (Section~\ref{sec:gravity-scaled}). The honest
population view is therefore not a froth but a near-empty field with a sparse handful
of round islands, each a lineage. The packed polygonal foam survives only as the
\emph{contrary-regime} illustration: what the supraverse would look like had the bounce
instanton been near-Planckian. We keep that as a stated hypothetical, not a figure.

\subsection{Time evolution within a bubble}

Show a single universe evolving from bounce (low-entropy boundary) through
cosmological evolution (expansion, structure formation, eventual heat death).
Highlight that time runs in a specific direction for each bubble, set by its
bounce.

\subsection{Cartasis membrane structure}

A detailed visualization of a single bounce membrane: parent matter falling in,
torsion building up, reaching critical density, bouncing through to the baby.
Show the flow of conserved quantities (energy, spin, chirality) across the
membrane.

\section{Implementation}

\paragraph{Tools.} Three.js or WebGL for browser-based interactive 3D
visualization; Python with Mayavi or Julia with Makie for high-quality scientific
plots; custom shader code for conformal coordinate transformations; and VR
(optional) for immersive supraverse exploration.

\paragraph{Data pipeline.} Physics-simulation outputs
(Chapter~\ref{chap:sims}) produce data describing supraverse structure; a
conformal transformation computes visualization coordinates; geometric primitives
(spheres for universes, surfaces for bounces, gradients for matter distributions)
are instantiated in a visualization scene; and interactive rendering provides
zoom, rotation, and color controls.

\paragraph{Aesthetic choices.} A print-safe grayscale scheme reads cleanly and is
the default (as in \Cref{fig:dilute} and the cover): a gray void, gray universe
interiors, and all contrast carried by the membranes---black for matter, white for
antimatter, with OG universes given a bold black-and-white double periphery. A colour
variant (matter in warm yellow/orange, antimatter in cool blue/purple) remains
available for screen use. In either scheme, bounce membranes are the prominent surfaces;
vacuum regions are subtle and low-contrast, possibly with faint
quantum-fluctuation speckle;
lighting that suggests ``gravity'' through subtle radial gradients on each
bubble; and, in animations, universes evolving forward in their internal time
while the supraverse as a whole appears static.

\section{``Art'' applications}

The visualization is genuinely beautiful as art independent of its scientific
use: cover art for the eventual paper/book, public-outreach materials, conference
presentations, educational tools, and generative art pieces inspired by
supraverse structure. This is a side benefit, not the primary purpose, but worth
noting because the framework has unusually visualizable structure compared to
most theoretical cosmology.

The cover of this book is one such piece: the dilute supraverse
(\Cref{fig:dilute}), gravity-scaled, set in portrait. We have deliberately
\emph{not} included the visually richer packed-polygonal ``lava lamp'' render here,
even as art, because it depicts the percolating regime the computation excludes
(Section~\ref{sec:instanton}); showing it as a centrepiece would misrepresent the
framework's own result. The honest centrepiece is the dilute scatter---sparser, but
true.

\subsection{What ``gravity-scaled'' means, computed---and the two-tier wallpaper}
\label{sec:gravity-scaled}

It pays to make the conformal map precise, because doing so changes the wallpaper's
verdict in an instructive way. Gravity-scaling renders objects of vastly different
physical size but comparable gravitational weight at comparable visual size by
rescaling each as if it had a common reference density $\rho_{\mathrm{ref}}$ (water,
say): a mass $M$ gets a gravity-scaled radius $R_{\mathrm{gs}}=(3M/4\pi
\rho_{\mathrm{ref}})^{1/3}$, magnified relative to its true size by
$\mu=(\rho/\rho_{\mathrm{ref}})^{1/3}$. The reference density cancels in any \emph{ratio}
of scaled radii, so the map is well-defined---only relative densities matter
(\Cref{fig:gravityscale}).

The result is two-tier, and the boundary is exactly the identity that defines a
universe. An OGU is a horizon-mass black hole, so its \emph{mean} density over its
Schwarzschild volume equals the cosmic critical density---which is $\sim$the void's
dark-energy density (the $R_s=R_H$ relation, $\Omega\approx1$, of Phase~0). Hence
$\rho_{\mathrm{OGU}}/\rho_{\mathrm{void}}\sim1$ and the OGU's gravity-scaling
magnification is $\mu\simeq1$: \emph{an OGU is not magnified relative to its void.}
At the OGU/void level, gravity-scaled coordinates \emph{equal} true coordinates, the
dilute scatter survives the scaling intact, and the wallpaper is honestly near-empty
in either system---OGUs remain tiny, round, and $\sim e^{I/4}$ horizons apart. One
level down it is the opposite: an astrophysical black hole is $\sim10^{17}\,
\mathrm{kg\,m^{-3}}$, denser than its host universe by factors that make
$\mu\sim10^{9}$--$10^{14}$, so gravity-scaling blows these otherwise-invisible compact
objects up until the nested next-generation structure becomes visible. \emph{That} is
where the conformal map earns its name. So the honest gravity-scaled supraverse is a
near-empty dilute scatter of OGUs, each a richly magnified nest within---and the
two-tier behaviour is itself a consequence of ``a universe is its own Schwarzschild
radius.'' Reproducible as \texttt{figures/scripts/ch10\_gravity\_scale.py}.

\paragraph{One number for the emptiness.} The cleanest single statement of how dilute
the supraverse is, the ratio of the mean OGU separation to an OGU's diameter. The
separation is $\sim e^{I/4}$ de Sitter horizon radii and an OGU (a horizon-mass hole)
is $\sim2$ horizon radii across, so
\[
  \frac{\text{distance between OGUs}}{\text{OGU diameter}}
  \;\sim\; \frac{e^{I/4}}{2}
  \;\sim\; 10^{\,2.7\times10^{83}} .
\]
This is the same gravity-scaled or not---an OGU and its void share a density, so
neither is magnified relative to the other. The \emph{exponent} alone,
$\sim2.7\times10^{83}$, already exceeds every number that appears in physics (the
observable universe holds $\sim10^{80}$ atoms): each universe is a lone speck whose
nearest sibling lies a gap away with more zeros than a universe has atoms. That is the
dilute verdict in one figure.

\begin{figure}[t]
  \centering
  \includegraphics[width=\textwidth]{gravity_scale}
  \caption{Gravity-scaling magnification $\mu=(\rho/\rho_{\mathrm{ref}})^{1/3}$.
    \emph{Left:} an OGU (horizon mass, mean density $=$ cosmic) sits at $\mu\simeq1$---
    gravity-scaled coordinates equal true coordinates at the OGU/void level---while
    astrophysical black holes, far denser than their host universe, magnify by
    $10^{9}$--$10^{14}$. \emph{Right:} the black-hole mean density ($\propto M^{-2}$)
    crosses the cosmic/void line precisely at the horizon mass, which is \emph{why} an
    OGU scales to unity (the $R_s=R_H$ identity).}
  \label{fig:gravityscale}
\end{figure}

\Cref{fig:dilute} is the honest counterpart: the dilute supraverse the instanton
computation selects (Chapter~\ref{chap:supraverse}, Section~\ref{sec:instanton}),
shown gravity-scaled so it can be seen at all. Because the vacuum has weight
($\rho_\Lambda\sim6\times10^{-27}\,\mathrm{kg\,m^{-3}}$), the void occupies real
volume, so universes are \emph{round} and isolated rather than packed polygons. The
crucial caveat is scale: in the dilute regime neighbouring OGUs sit
$\sim e^{I/4}$ horizons apart ($I\sim10^{84}$), so a \emph{linear} ``true to scale''
frame would be uniform void with at most a single infinitesimal dot---the dilute
verdict made visible, but un-renderable as an image. The figure is therefore
gravity-scaled and illustrative in its sizes and spacings; what is physical is the
qualitative content (isolated, round, well-separated discs in a dominating weighted
void, never a packed polygonal foam). The packed render above is the legible
percolating limit; neither is a linear map.

\begin{figure}[t]
  \centering
  \includegraphics[width=\textwidth]{dilute_supraverse}
  \caption{The dilute supraverse, gravity-scaled (sizes and spacings illustrative, not
    linear). The weighted void ($\rho_\Lambda\sim6\times10^{-27}\,\mathrm{kg\,m^{-3}}$,
    faintly speckled with zero-point energy) dominates; universes are sparse,
    \emph{round}, isolated islands (gravity wells with chirality-toned Cartasis
    membranes, faint nested $\mathrm{BHU}$ children). The inset zooms one, marking
    \emph{home}. At a literal linear scale the frame would be near-empty---neighbours
    lie $\sim e^{I/4}$ horizons apart---so the void's weight keeps the islands round
    while the dilution makes the true gaps astronomically larger than drawn.}
  \label{fig:dilute}
\end{figure}

\section{Computational notes}

\paragraph{Static vs.\ dynamic.} The supraverse is static (a block universe), but
each bubble has its own internal dynamics. The visualization can be a static
snapshot, an animation showing each bubble evolving in its own time while
supraverse-level structure stays fixed, or interactive exploration through the
foam.

\paragraph{Scale handling.} Universes span many orders of magnitude. The
conformal scaling helps, but additional logarithmic mapping might be needed for
some visualizations.

\paragraph{Performance.} Real-time interactive exploration with many nested
bubbles is computationally demanding. Optimizations: level-of-detail rendering,
bubble culling, and asynchronous loading of detailed sub-structures.

\section{Specific plot ideas}

\begin{enumerate}
  \item \textbf{``The Cartasis membrane.''} A 2D cross-section showing the
    density profile of matter falling into a black hole, bouncing at $\rhoc$, and
    the geometric transition into the baby universe.
  \item \textbf{``Tardis effect.''} Side-by-side external view (small black-hole
    sphere) and internal view (vast cosmological space), annotated to show they
    are connected through the membrane.
  \item \textbf{``Lineage tree.''} A family-tree visualization with our universe
    at the bottom, ancestors going up, each generation alternating chirality
    colors, and the OG ancestor at the top.
  \item \textbf{``Supraverse foam.''} A wide-field view of many independent
    lineages distributed across cosmic vacuum, with the gravity-scaled mapping
    making all universes appear comparable in size.
  \item \textbf{``Time arrow.''} A diagram showing how time flows away from each
    bounce in each universe, with different arrows pointing different directions
    in supraverse coordinates.
  \item \textbf{``CMB $=$ Hawking radiation.''} A diagram showing the geometric
    identification of the parent's event horizon with the baby's cosmic horizon.
  \item \textbf{``Dark matter as parent matter.''} A diagram showing the parent's
    matter distribution projected through the Cartasis membrane, manifesting as
    dark-matter halos in the baby.
\end{enumerate}

\section{First implementation: simple Python prototypes}

To get started without major infrastructure:

\begin{verbatim}
import matplotlib.pyplot as plt
import numpy as np

def conformal_scale(rho):
    """Gravity-based scaling factor"""
    return (G * rho)**(-0.5)

def render_universe(ax, x, y, mass, chirality):
    """Render a single universe bubble"""
    radius = conformal_scale(mass / volume(mass))
    color = 'orange' if chirality > 0 else 'blue'
    circle = plt.Circle((x, y), radius, color=color, alpha=0.5)
    ax.add_patch(circle)

def render_supraverse(universes):
    """Render a collection of universes"""
    fig, ax = plt.subplots()
    for u in universes:
        render_universe(ax, u.x, u.y, u.mass, u.chirality)
    plt.show()
\end{verbatim}

This is the simplest possible starting point. Real visualizations would be much
more sophisticated, but the principle is the same: use gravity-based scaling to
make multi-scale supraverse structure visible.
